% Generated draft: requires mathematical and source review.
\documentclass[11pt]{article}
\usepackage[margin=25mm]{geometry}
\usepackage{fontspec}
\setmainfont{Latin Modern Roman}
\usepackage{amsmath,amssymb,mathtools,unicode-math}
\setmathfont{Latin Modern Math}
\usepackage{enumitem,needspace,xurl,hyperref}
\setlength{\parindent}{0pt}
\setlength{\parskip}{5pt}
\setlength{\emergencystretch}{3em}
\newcommand{\N}{\mathbb{N}}
\newcommand{\Z}{\mathbb{Z}}
\newcommand{\Q}{\mathbb{Q}}
\newcommand{\R}{\mathbb{R}}
\newcommand{\C}{\mathbb{C}}
\newcommand{\F}{\mathbb{F}}
\newcommand{\A}{\mathbb{A}}
\newcommand{\PP}{\mathbb{P}}
\newcommand{\E}{\mathbb{E}}
\newcommand{\eps}{\varepsilon}
\newcommand{\dd}{\,\mathrm{d}}
\begin{document}
\section*{UnsolvedMath ID: 30004549}
\textbf{Source problem code:} catalogue entry
\subsection{Definitions and mathematical statement}
Let \(M\) be a compact real \(4\)-manifold admitting an almost-complex structure. Using the conventional smooth interpretation, an almost-complex structure is a smooth bundle endomorphism \(J:TM\to TM\) satisfying \(J^2=-\operatorname{id}_{TM}\). Write \(\mathcal J(M)\) for the set of such structures. Using real smooth differential forms, define
\[
\mathcal Z_J^-:=\left\{\alpha\in\Omega^2(M;\R):\ \mathrm d\alpha=0,\quad \alpha(JX,JY)=-\alpha(X,Y)\text{ for all smooth vector fields }X,Y\right\},
\qquad h_J^-:=\dim_{\R}\mathcal Z_J^-.
\]
Here integrability means that \(J\) is induced by a complex manifold structure on \(M\).

The supplied record poses two assertions, intended for every compact real \(4\)-manifold \(M\) admitting an almost-complex structure:
\begin{enumerate}
\item The set
\[
\mathcal G(M):=\{J\in\mathcal J(M):h_J^-=0\}
\]
is generic in \(\mathcal J(M)\).
\item For every \(J\in\mathcal J(M)\),
\[
h_J^-\geq 3\quad\Longrightarrow\quad J\text{ is integrable}.
\]
\end{enumerate}
The source does not specify a topology on \(\mathcal J(M)\) or whether ``generic'' means residual, open and dense, or another condition; consequently the first assertion is not fully specified. Smoothness and real coefficients above are explicit conventional interpretations requiring source confirmation. The record also does not state whether \(M\) must be without boundary. The displayed integrability definition has its usual complex-manifold meaning when \(M\) has no boundary; any intended boundary case requires additional conventions. No finite-dimensionality assertion is made without resolving these hypotheses. The definition of \(h_J^-\) follows the supplied statement's space of closed forms; the title's reference to cohomology does not by itself justify substituting a cohomological definition.

The supplied dated curation labels the problem partially solved: it reports generic vanishing as proved, while describing the implication \(h_J^-\geq3\Rightarrow J\text{ integrable}\) as conjectural, with only bounds and special cases known. These are supplied literature assessments, not independently verified conclusions.
\subsection{Short English statement}
The record asks whether a generic almost-complex structure on a compact real four-manifold has no nonzero closed anti-invariant real two-forms, and whether the existence of at least three linearly independent such forms forces integrability. Supplied dated curation reports the first assertion proved and the second conjectural. The meaning of genericity and the intended boundary and regularity conventions require confirmation.
\subsection{Sources}
\begin{itemize}
\item Supplied local UnsolvedMath registry record, permanent ID \(30004549\), source code \texttt{catalogue entry}, titled ``Generic Vanishing of Anti-Invariant Cohomology in Dimension Four,'' in the Oberwolfach Reports Open Problems collection. No individual proposer is supplied; the proposed year is recorded as 2020.
\item Supplied source citation: \emph{Mini-Workshop: Almost Complex Geometry (2021)}, Oberwolfach Reports, pp.\ 1657--1691, DOI \texttt{10.4171/owr/2020/33}. \url{https://doi.org/10.4171/owr/2020/33}. The background attributes the problem to workshop open-problem discussions documented by Oberwolfach Reports.
\item Supplied literature reference: Tan--Wang--Zhang, 2015. \url{https://doi.org/10.1007/s12220-014-9477-2}. No article title or more detailed bibliographic information is supplied.
\item Supplied literature reference: Oberwolfach Report 33/2020. \url{https://doi.org/10.4171/owr/2020/33}.
\item Literature assessment checked at \texttt{2026-08-22}, according to the supplied record: generic vanishing is reported proved, while the integrability implication remains conjectural, with bounds and special cases known. This is supplied dated curation derived from the OWR record, not an independently verified assessment; time-sensitive status and cited claims require independent verification.
\end{itemize}
\end{document}
